\par\noindent\textbf{(1)} 相邻两弹簧给m1ü\_n=K(v\_n+v\_\{n−1\}−2u\_n)、m2v̈\_n=K(u\_n+u\_\{n+1\}−2v\_n)。整链同向平移时各弹簧伸长为零，保证声学模在q=0必有零频Goldstone分支。
\par\noindent\textbf{(2)} 设un=Uexp[i(qna−ωt)]、vn=Vexp[i(qna−ωt)]，动力矩阵行列式为(2K−m1ω²)(2K−m2ω²)−4K²cos²(qa/2)=0，即m1m2ω⁴−2K(m1+m2)ω²+4K²sin²(qa/2)=0。
\par\noindent\textbf{(3)} 解出ω²\_±=K(1/m1+1/m2)±K√[(1/m1+1/m2)²−4sin²(qa/2)/(m1m2)]。在q=0，低分支ω\_−=0，高分支ω\_+²=2K(1/m1+1/m2)，分别代表整体平移声学模与胞内反向振动光学模。
\par\noindent\textbf{(4)} 区边界qa=π时ω\_−²=2K/max(m1,m2)、ω\_+²=2K/min(m1,m2)。m1≠m2时出现有限禁频区；m1=m2时两支恰在折叠区边界相接，反映人为选择双原子晶胞产生的带折叠。
